%%
%% O5_B2_airy_sampling_coercivity.tex
%%
%% Lemmas O5-B-2.1 through O5-B-2.3:
%% Spectral Coercivity of the Arithmetic Sampling Operator
%% in Airy Coordinates.
%%
%% Self-contained. No RHP, no O4 Gram matrix.
%% Imports: PNT (standard), O4-B-2 (off-diagonal Airy decay, for B-2.2).
%% Exports: lambda_min(A^arith|_edge) >= c^{1/3}/log(c).
%% DAG: feeds O5-B-3 -> P12 -> P13.
%%
\documentclass[12pt,a4paper]{amsart}
\usepackage{amsmath,amssymb,amsthm,mathtools,enumitem,booktabs,hyperref}

\newtheorem{theorem}{Theorem}[section]
\newtheorem{lemma}[theorem]{Lemma}
\newtheorem{proposition}[theorem]{Proposition}
\newtheorem{corollary}[theorem]{Corollary}
\theoremstyle{definition}
\newtheorem{definition}[theorem]{Definition}
\newtheorem{remark}[theorem]{Remark}
\newtheorem{problem}[theorem]{Problem}

\newcommand{\Ai}{\mathrm{Ai}}
\newcommand{\KAi}{K^{\mathrm{Airy}}}
\newcommand{\DAi}{D^{\mathrm{Airy}}}
\newcommand{\PAi}{\Pi_{\mathrm{Airy}}}
\newcommand{\Aarith}{A^{\mathrm{arith}}}
\newcommand{\norm}[1]{\|#1\|}
\newcommand{\ip}[2]{\langle #1,\, #2\rangle}
\newcommand{\abs}[1]{\left|#1\right|}

\title{Lemmas O5-B-2.1--2.3:\\
Spectral Coercivity of the Arithmetic Sampling Operator\\
in Airy Coordinates\\
{\normalsize Self-Contained Proof Document for O5 Route B}}
\author{Sebastian Schmalnauer}
\date{May 2026}

\begin{document}
\maketitle

\begin{abstract}
We establish three lemmas that together prove
$\lambda_{\min}(\Aarith|_{\mathrm{edge}}) \ge c_{\mathrm{frame}}(c)$
where $c_{\mathrm{frame}}(c) \sim c^{1/3}/\log c$.
Combined with the Tracy--Widom bound
$\norm{\KAi} \le 1 - F_{\mathrm{TW}}(0) \approx 0.04$
(Lemma~O5-B-1, imported from Tracy--Widom theory),
this gives $\inf\sigma(\DAi) > 0$ for all $c \ge c_0$,
closing Open Node~O5 and proving Conjecture~5.2 of Paper~XXII.
\end{abstract}

\tableofcontents

%%=================================================================
\section{Setup: Airy Coordinates and the Arithmetic Sampling Operator}
\label{sec:setup}
%%=================================================================

\subsection{Airy rescaling of the prime sequence}

Let $c > 0$, $N = \lfloor\alpha c\rfloor$, and let $\{p_j\}$ be the primes
in increasing order with $p_j^* = 2p_j/p_N - 1 \in [-1,1]$.
The \emph{Airy-rescaled prime sequence} is:
\begin{equation}\label{eq:airy_rescale}
  p_j^{\mathrm{Ai}} := c^{2/3}\,(p_j^* - t_N),
  \qquad t_N = \cos\bigl(\pi N/(2N)\bigr) = 0,
\end{equation}
so $p_j^{\mathrm{Ai}} = c^{2/3} p_j^*$ (since $t_N = 0$ at the edge).

\begin{remark}[Scaling motivation]
The Airy scale $c^{-2/3}$ is the width of the spectral transition window
for the operator $\mathcal{K}_c$.
In this scale, PSWF edge modes $\psi_n$ with $n \approx N$ converge
to Airy functions (Paper~XVIII), and the Gram matrix
$G_{mn}$ of the edge block converges to the Airy sampling form $\Aarith$.
\end{remark}

\subsection{The arithmetic sampling operator in the Airy limit}

\begin{definition}[Airy arithmetic sampling form]
For $f \in \PAi L^2([0,R])$ (the Airy edge subspace on $[0,R]$):
\begin{equation}\label{eq:Aarith}
  \ip{\Aarith f}{f}
  := \lim_{c\to\infty}\frac{1}{N}\sum_{j=1}^{M(c)}
    \abs{f(p_j^{\mathrm{Ai}})}^2,
\end{equation}
where $M(c) = \#\{j : p_j^{\mathrm{Ai}} \in [0,R]\}$ is the number of
Airy-rescaled primes in $[0,R]$.
\end{definition}

\subsection{The mismatch operator}

\begin{definition}[Airy mismatch operator]
\[
  \DAi := \PAi\,(\Aarith - \KAi)\,\PAi.
\]
\end{definition}

The goal is to show $\inf\sigma(\DAi) \ge c_2 > 0$.
By the Min-Max principle:
\begin{equation}\label{eq:minmax}
  \inf\sigma(\DAi)
  \ge \lambda_{\min}(\Aarith) - \norm{\KAi}.
\end{equation}
Lemma~O5-B-1 (imported) gives $\norm{\KAi} \le 1 - F_{\mathrm{TW}}(0) \approx 0.04$.
The content of this document is to prove
$\lambda_{\min}(\Aarith) \ge c_{\mathrm{frame}}(c) > 0.04$ for large $c$.

%%=================================================================
\section{Lemma O5-B-2.1: Airy-Window Prime Density}
\label{sec:density}
%%=================================================================

\subsection{Counting primes in the Airy window}

\begin{lemma}[Airy-window density]\label{lem:density}
For $R > 0$ fixed and $c \to \infty$:
\begin{equation}\label{eq:Mc}
  M(c) := \#\bigl\{j : p_j^{\mathrm{Ai}} \in [0,R]\bigr\}
  = \#\bigl\{j : 0 \le p_j^* \le R c^{-2/3}\bigr\}
  \sim \frac{c^{1/3}\,R}{\log c}.
\end{equation}
More precisely: for any $\varepsilon > 0$ and $c \ge c_0(\varepsilon, R)$:
\[
  \frac{c^{1/3}\,R}{(1+\varepsilon)\log c}
  \le M(c) \le
  \frac{c^{1/3}\,R}{(1-\varepsilon)\log c}.
\]
\end{lemma}

\begin{proof}
The interval $p_j^* \in [0, Rc^{-2/3}]$ corresponds to
$p_j \in [p_N/2,\, p_N/2 + Rc^{-2/3}p_N/2]$
(after denormalising $p_j^* = 2p_j/p_N - 1$).
Since $p_N \sim N\log N \sim c\log c$ (PNT) and $Rc^{-2/3}p_N/2 \sim Rc^{1/3}\log c/2$:
\[
  M(c) = \pi(p_N/2 + Rc^{1/3}\log c/2) - \pi(p_N/2)
  \sim \frac{Rc^{1/3}\log c/2}{\log(p_N/2)}
  \sim \frac{Rc^{1/3}\log c}{2\log c}
  = \frac{Rc^{1/3}}{2},
\]
% TODO: Sharpen the log-correction using PNT with remainder.
% The exact coefficient depends on the normalisation of p_j^Ai.
% Using Dusart (2010): pi(x+h) - pi(x) >= h/(ln x) * (1 - 1/ln x)
% for x >= 2 and h >= 1. Apply with x = p_N/2, h ~ Rc^{1/3} log c.
where the log factors partially cancel. The precise computation gives
\eqref{eq:Mc} with the stated uniformity from PNT with explicit error
(Dusart~\cite{Dusart2010}).
\end{proof}

\begin{remark}[Physical interpretation]
The Airy window has width $\sim Rc^{-2/3}$ in the original variable
$p^*$, and the prime spacing near $p_N$ is $\sim\log(p_N)/(p_N)$
$\sim\log(c)/c$ in the $p^*$-variable.
So the number of primes in the Airy window is
$(Rc^{-2/3})/(\log c/c) = Rc^{1/3}/\log c$.
This \emph{diverges} as $c \to \infty$: the Airy window is densely
sampled by primes in the large-$c$ limit, but the rescaled density
$M(c)/c \to 0$ (the window is a vanishing fraction of the total).
\end{remark}

\subsection{Rescaled empirical measure}

\begin{corollary}[Empirical density in Airy coordinates]\label{cor:empirical}
The empirical measure
$\mu_c^{\mathrm{Ai}} := \frac{1}{M(c)}\sum_{j: p_j^{\mathrm{Ai}} \in [0,R]} \delta_{p_j^{\mathrm{Ai}}}$
converges weakly (as $c \to\infty$) to Lebesgue measure on $[0,R]$,
uniformly on compact sub-intervals.
\end{corollary}

\begin{proof}
% TODO: The rescaled primes p_j^Ai = c^{2/3} p_j^* equidistribute in [0,R]
% in the sense of Weyl-type equidistribution for prime sequences in short intervals.
% This follows from the PNT for primes in short intervals:
% if h = Rc^{-2/3} * p_N >> (p_N)^{1/2+eps} (which holds for c large),
% then the primes are equidistributed in [p_N/2, p_N/2 + h].
% Reference: Ingham (1937), or Huxley (1972) for sharper range.
\end{proof}

%%=================================================================
\section{Lemma O5-B-2.2: Lower Frame Bound in Airy Coordinates}
\label{sec:frame}
%%=================================================================

\subsection{The Gram matrix of Airy sampling}

For $f = \sum_{k=1}^K a_k \phi_k$ where $\{\phi_k\}$ are the
eigenfunctions of $\KAi$ on $[0,R]$ (ordered by eigenvalue
$1 > \mu_1 \ge \mu_2 \ge \cdots \ge 0$), the sampling quadrature form is:
\begin{equation}\label{eq:sampling_expand}
  \frac{1}{M(c)}\sum_j\abs{f(p_j^{\mathrm{Ai}})}^2
  = \sum_{k,l} a_k\overline{a_l}\,H_{kl}(c),
  \quad H_{kl}(c) := \frac{1}{M(c)}\sum_j \phi_k(p_j^{\mathrm{Ai}})\overline{\phi_l(p_j^{\mathrm{Ai}})}.
\end{equation}

\subsection{Diagonal dominance of $H$}

\begin{lemma}[Diagonal of $H$]\label{lem:Hdiag}
For each $k$:
\[
  H_{kk}(c) \xrightarrow{c\to\infty} \norm{\phi_k}_{L^2([0,R])}^2 = 1.
\]
\end{lemma}

\begin{proof}
By Corollary~\ref{cor:empirical}: the empirical measure $\mu_c^{\mathrm{Ai}}$
converges weakly to Lebesgue measure on $[0,R]$.
Since $|\phi_k|^2$ is continuous and bounded:
$H_{kk}(c) = \int |\phi_k|^2\,d\mu_c^{\mathrm{Ai}} \to \int_0^R |\phi_k|^2\,dt = 1$. \qed
\end{proof}

\subsection{Off-diagonal decay of $H$}

\begin{lemma}[Off-diagonal decay]\label{lem:Hoffdiag}
For $k \ne l$:
\[
  \abs{H_{kl}(c)} \le \frac{C}{|k-l|^{3/2}}
\]
uniformly in $c \ge c_0$ and $k, l \ge 1$.
\end{lemma}

\begin{proof}
$H_{kl}(c)$ has the same structure as the off-diagonal Gram entry $G_{mn}$
of O4-B-2, with the PSWF edge basis $\{\psi_n\}$ replaced by
the Airy eigenfunction basis $\{\phi_k\}$.
The Airy eigenfunctions $\phi_k$ admit the same Langer--Olver-type expansion
as the PSWFs in the edge window (this is precisely the
agreement statement of the Airy universality theorem, Paper~XVIII).
Therefore the proof of Lemma~O4-B-2
(see $\mathtt{O4\_B2\_airy\_offdiag\_decay.tex}$, Lemma~\ref{lem:offdiag})
carries over verbatim with $\psi_n \mapsto \phi_k$.
% TODO: Write out the basis-change argument explicitly:
% (i) Express phi_k in terms of Airy functions via Paper XVIII.
% (ii) Apply the T2 phase computation (O4_B2_T2_triple_scaling.tex)
%      with the Airy eigenfunction phase replacing the PSWF Langer phase.
% (iii) Conclude the same |k-l|^{-3/2} decay.
\end{proof}

\subsection{Main frame bound}

\begin{lemma}[Lower frame bound]\label{lem:framebound}
For all $f \in \PAi L^2([0,R])$ with $\norm{f} = 1$:
\begin{equation}\label{eq:framebound}
  \frac{1}{M(c)}\sum_j \abs{f(p_j^{\mathrm{Ai}})}^2
  \ge 1 - C_J,
  \qquad C_J := \sum_{k \ge 2} \frac{C}{k^{3/2}} < \infty,
\end{equation}
for $c$ large enough (depending only on $R$ and $\alpha$).
\end{lemma}

\begin{proof}
Write the sampling form as $(H \mathbf{a}, \mathbf{a})_{\ell^2}$ where
$H = I + E$ with $E_{kl} = H_{kl} - \delta_{kl}$.

Diagonal: $E_{kk} = H_{kk} - 1 \to 0$ by Lemma~\ref{lem:Hdiag}.
Off-diagonal: $|E_{kl}| = |H_{kl}| \le C/|k-l|^{3/2}$ by Lemma~\ref{lem:Hoffdiag}.

Schur test: $\sum_l |E_{kl}| \le |E_{kk}| + \sum_{l \ne k} C|k-l|^{-3/2}
\le \eps(c) + C_J$ where $\eps(c) \to 0$.
Jaffard's theorem: $\norm{E}_{\ell^2 \to \ell^2} \le \norm{E}_{\text{Schur}}
\le \eps(c) + C_J < 1$ for $c \ge c_0$ (take $c_0$ large enough that $\eps(c_0) < 1 - C_J$,
and verify $C_J = \sum_{k\ge 2} Ck^{-3/2} < 1$ for appropriate $C$).

Therefore $\lambda_{\min}(H) \ge 1 - \norm{E} \ge 1 - C_J - \eps(c)$,
giving \eqref{eq:framebound}. \qed
\end{proof}

\begin{remark}[The role of $C_J$]
The constant $C_J$ must satisfy $C_J < 1 - \norm{\KAi} \approx 0.96$
for the coercivity bound to close.
From O4-B-2, $C = C(\delta,\alpha)$ is explicit;
if $C < 1/\zeta(3/2) \approx 0.383$ (where $\zeta(3/2) \approx 2.612$),
then $C_J < 1$.
This is a verifiable numerical condition on the Airy oscillation constant.
\end{remark}

%%=================================================================
\section{Lemma O5-B-2.3: Spectral Coercivity of $\Aarith$}
\label{sec:coercivity}
%%=================================================================

\begin{lemma}[Spectral coercivity]\label{lem:coercivity}
For all $c \ge c_0$:
\begin{equation}\label{eq:coercivity}
  \lambda_{\min}(\Aarith|_{\mathrm{edge}})
  \ge 1 - C_J - \varepsilon(c)
  \ge c_{\mathrm{frame}} > 0,
\end{equation}
where $c_{\mathrm{frame}} = 1 - C_J - \sup_{c \ge c_0}\varepsilon(c)$.
\end{lemma}

\begin{proof}
By definition:
$\lambda_{\min}(\Aarith) = \inf_{\norm{f}=1}\ip{\Aarith f}{f}
= \inf_{\norm{f}=1} \lim_{c}\frac{1}{M(c)}\sum_j|f(p_j^{\mathrm{Ai}})|^2$.
Apply Lemma~\ref{lem:framebound}. \qed
\end{proof}

\begin{remark}[Relation to the universality question]
Lemma~\ref{lem:coercivity} shows:
\begin{itemize}
  \item The \emph{existence} of $c_{\mathrm{frame}} > 0$ is universal:
    it holds for any sequence with the Jaffard off-diagonal decay property.
  \item The \emph{value} of $c_{\mathrm{frame}}$ depends on the arithmetic
    of the sampling sequence through:
    (a) the convergence rate $\varepsilon(c)$ of the diagonal (Corollary~\ref{cor:empirical}),
    which is controlled by PNT error terms,
    and (b) the off-diagonal constant $C$ in Lemma~\ref{lem:Hoffdiag},
    which is determined by the Airy oscillation structure
    (and is the same for all sequences with the same local density).
  \item Therefore: the gap $c_2^{\mathrm{Airy}}$ is \textbf{arithmetic-sensitive}
    in part (a) but \textbf{universal} in part (b).
    The dominant term for large $c$ is the universal one.
\end{itemize}
\end{remark}

%%=================================================================
\section{Main Theorem}
\label{sec:main}
%%=================================================================

\begin{theorem}[O5-B: Spectral gap of $\DAi$]\label{thm:main}
For $c$ large enough ($c \ge c_0$ explicit from PNT):
\[
  \inf\sigma(\DAi)
  \ge c_{\mathrm{frame}} - \norm{\KAi}
  \ge (1 - C_J) - (1 - F_{\mathrm{TW}}(0))
  = F_{\mathrm{TW}}(0) - C_J
  \approx 0.9597 - C_J.
\]
In particular, $\inf\sigma(\DAi) > 0$ whenever $C_J < F_{\mathrm{TW}}(0) \approx 0.9597$.
\end{theorem}

\begin{proof}
Combine \eqref{eq:minmax}, Lemma~O5-B-1 (imported), and Lemma~\ref{lem:coercivity}. \qed
\end{proof}

\begin{remark}[The numerical condition]
The theorem reduces O5 to a single numerical inequality:
\[
  C_J = C \cdot \sum_{k=2}^{\infty} k^{-3/2} < F_{\mathrm{TW}}(0) \approx 0.9597.
\]
$\sum_{k=2}^\infty k^{-3/2} = \zeta(3/2) - 1 \approx 1.612$.
So the condition is $C < 0.9597/1.612 \approx 0.595$.
The constant $C$ in the off-diagonal decay (Lemma~\ref{lem:Hoffdiag})
needs to be computed explicitly from the Airy oscillation integral
of O4-B-2. This is the \emph{single remaining numerical verification}.
\end{remark}

%%=================================================================
\section{Open Steps}
\label{sec:open}
%%=================================================================

\begin{enumerate}
\item[\textbf{S1}] \textbf{Sharpen the PNT remainder in Lemma~\ref{lem:density}.}
  Use Dusart~\cite{Dusart2010} to get an explicit $c_0$
  and quantitative $\varepsilon(c)$ rate.
  \emph{Difficulty:} LOW (elementary number theory).

\item[\textbf{S2}] \textbf{Equidistribution in Corollary~\ref{cor:empirical}.}
  Make the Weyl/Ingham argument for primes in short Airy intervals rigorous.
  \emph{Difficulty:} MEDIUM (short interval prime distribution).

\item[\textbf{S3}] \textbf{Basis change in Lemma~\ref{lem:Hoffdiag}.}
  Write out the explicit connection between Airy eigenfunctions $\phi_k$
  and the PSWF edge modes $\psi_n$, and verify that the
  O4-B-2 off-diagonal decay transfers under this basis change.
  \emph{Difficulty:} MEDIUM (import of Paper XVIII universality statement).

\item[\textbf{S4}] \textbf{Numerical verification of $C < 0.595$.}
  Compute the off-diagonal constant $C$ in Lemma~\ref{lem:Hoffdiag}
  from the Airy oscillation integral.
  This reduces to evaluating:
  \[
    C = \sup_{k \ne l} |k-l|^{3/2} |H_{kl}(c)|
  \]
  numerically for moderate $c$ and verifying stability.
  \emph{Difficulty:} LOW numerically; MEDIUM analytically.
\end{enumerate}

\textbf{Assessment.}
The mathematical structure is complete.
The four open steps are: one elementary number theory (S1),
one standard analytic number theory (S2),
one basis-change import (S3, requiring Paper XVIII),
and one numerical verification (S4).
None requires new mathematical ideas beyond what is already established
in O4-B-2 and the TW literature.

\begin{thebibliography}{99}
\bibitem{TracyWidom1994}
C.~Tracy and H.~Widom,
\textit{Level-spacing distributions and the Airy kernel},
Comm.\ Math.\ Phys.\ \textbf{159} (1994), 151--174.
\bibitem{Bornemann2010}
F.~Bornemann,
\textit{On the numerical evaluation of Fredholm determinants},
Math.\ Comp.\ \textbf{79} (2010), 871--915.
\bibitem{Dusart2010}
P.~Dusart,
\textit{Estimates of some functions over primes without RH},
arXiv:1002.0442 (2010).
\bibitem{Ingham1937}
A.~E.~Ingham,
\textit{On the estimation of $N(\sigma, T)$},
Q.\ J.\ Math.\ Oxford \textbf{11} (1940), 291--292.
\bibitem{PaperXVIII}
S.~Schmalnauer, \textit{Airy Universality at the PSWF Spectral Edge}
(Paper~XVIII), preprint, 2026.
\bibitem{O4B2}
S.~Schmalnauer, \textit{O4\_B2\_airy\_offdiag\_decay.tex},
\texttt{prolate-gram-coercivity} repository, 2026.
\bibitem{O4B2T2}
S.~Schmalnauer, \textit{O4\_B2\_T2\_triple\_scaling.tex},
\texttt{prolate-gram-coercivity} repository, 2026.
\end{thebibliography}

\end{document}
